%% CSE 578 Final Project. Individual Contributions Report.
%% Companion to report.tex. Two pages max. Compile with pdfLaTeX.

\documentclass[sigconf,nonacm,natbib=false]{acmart}

\settopmatter{printacmref=false}
\renewcommand\footnotetextcopyrightpermission[1]{}
\pagestyle{plain}

\usepackage{booktabs}
\usepackage{hyperref}
\usepackage{xcolor}
\usepackage{microtype}

\begin{document}

\title{Individual Contributions Report}
\subtitle{Color is Destiny, or is it? CSE 578 Final Project, Team 7}

\author{Mathew Abraham Palliambil}
\affiliation{\institution{Arizona State University}}
\email{mpalliam@asu.edu}

\author{Asmit Datta}
\affiliation{\institution{Arizona State University}}
\email{adatta18@asu.edu}

\author{Basil Shibu Thomas}
\affiliation{\institution{Arizona State University}}
\email{bthoma70@asu.edu}

\author{Kruthika Kadurhalli Raghu}
\affiliation{\institution{Arizona State University}}
\email{kkadurha@asu.edu}

\author{Chandana Priya Srinivasa}
\affiliation{\institution{Arizona State University}}
\email{csrini10@asu.edu}

\author{Parth Jain}
\affiliation{\institution{Arizona State University}}
\email{pjain121@asu.edu}

\renewcommand{\shortauthors}{Team 7}

\maketitle

%% ============================================================
\section{Team Split}
%% ============================================================

The six member team divided into two sub-teams of three. Side~A
(\emph{Color is Destiny}, charts A1 to A6) was carried by Mathew,
Asmit, and Basil. Side~B (\emph{Color is Irrelevant}, charts B1 to
B6) was carried by Kruthika, Chandana, and Parth. Three members
held additional cross-cutting roles: Mathew led project
coordination, Asmit led web development and chart implementation,
and Basil led visual and poster design.

For each owned chart, the responsible member proposed the
analytical question, specified the SQL aggregation, drafted the
on-page narrative copy, and validated the rendered chart against
the intended argument. Asmit then compiled each member's chart
work into a single scrollytelling webpage.

\begin{table}[h]
\centering
\small
\begin{tabular}{@{}lll@{}}
\toprule
\textbf{Member} & \textbf{Charts} & \textbf{Cross-cutting role} \\
\midrule
Mathew A.\ Palliambil  & A2, A4 & Project lead \\
Asmit Datta            & A1, A6 & Lead developer, web integration \\
Basil Shibu Thomas     & A3, A5 & Poster and visual design \\
Kruthika K.\ Raghu     & B1, B5 & Side B analysis \\
Chandana P.\ Srinivasa & B2, B3 & Side B analysis \\
Parth Jain             & B4, B6 & Side B analysis \\
\bottomrule
\end{tabular}
\caption{Sub-team split and roles.}
\label{tab:roles}
\end{table}

%% ============================================================
\section{Mathew Abraham Palliambil}
%% ============================================================
\textbf{Email:} \texttt{mpalliam@asu.edu}.
\textbf{Role:} project lead.

Mathew owned the project plan from kickoff to submission. He
decomposed the rubric into deliverables (data, pipeline, narrative,
12 charts, two innovative views, scrollytelling site, poster,
report), assigned ownership, and maintained the running checklist
that anchored every check-in. When the team pivoted twice
(Aggressive vs.\ Defensive, then Theory vs.\ Weird, then Color is
Destiny), he rebuilt the task plan and reallocated chart ownership
so the schedule absorbed the pivot.

\textbf{Charts.}
A2 (White edge across Elo bands): designed the question and the
per-band aggregation.
A4 (Top White-winning openings, treemap): set the volume threshold
($\geq 2{,}000$ games) and authored the ``deep bench'' framing.

%% ============================================================
\section{Asmit Datta}
%% ============================================================
\textbf{Email:} \texttt{adatta18@asu.edu}.
\textbf{Role:} lead developer and web integration.

Asmit compiled each team member's chart work into a single
scrollytelling webpage. He took each owner's analytical
specification and narrative copy and integrated them into the live
site. Concretely, he wrote the DuckDB preprocessing script
(\texttt{preprocess\_color.py}) that reduces the 4.4~GB CSV to 13
small JSONs, built the static site
(\texttt{index.html}, \texttt{side\_a.html}, \texttt{side\_b.html}),
implemented the scrollytelling controller on
\texttt{IntersectionObserver}, and rendered the 12 D3.js v7 charts
from the per-chart specifications produced by their owners. He also
implemented the A3 to A6 cross-chart click-jump and set up the Git
repository and reproducibility instructions.

\textbf{Charts.}
A1 (Headline bar): defined the 50.4 / 3.1 / 46.5 split that anchors
the site.
A6 (Tug-of-War, innovative): co-designed the rope-and-knot encoding
across 13 categorical cuts on a single shared axis and implemented
its geometry in D3.

%% ============================================================
\section{Basil Shibu Thomas}
%% ============================================================
\textbf{Email:} \texttt{bthoma70@asu.edu}.
\textbf{Role:} poster and visual design.

Basil owned the visual identity of the project. He designed the
25\,$\times$\,30 inch final poster end to end: layout, typography,
color, the side-by-side Side~A and Side~B split, the headline
callouts (+3.9, 19, 61, 0.9, 0.4), and the team strip. He also
defined the shared palette and the aesthetic contrast between the
two scrollytelling sides (a near-white serif voice for Side~A and a
dark editorial voice for Side~B), so a viewer can tell which side
they are reading at a glance.

\textbf{Charts.}
A3 (Edge across time controls): designed the
bullet/blitz/rapid/classical comparison.
A5 (White-mate heatmap): chose the \texttt{interpolateYlOrRd} ramp
so heat reads against the page background, replacing an earlier
greyscale version that vanished.

%% ============================================================
\section{Kruthika Kadurhalli Raghu}
%% ============================================================
\textbf{Email:} \texttt{kkadurha@asu.edu}.
\textbf{Role:} Side~B analysis.

\textbf{Charts.}
B1 (Waffle, ``out of 100''): designed the reframing that converts
the same +3.9pp number Side~A opens with into Side~B's
philosophical pivot (4 cells color, 96 cells everything else).
B5 (Rating-gap stacked bars): specified the seven-bucket ladder on
\texttt{WhiteElo} minus \texttt{BlackElo} and surfaced the project's
killer comparison: a 61.2pp swing across rating gap versus 3.9pp
across color. B5 is the chart that quantitatively closes Side~B.

%% ============================================================
\section{Chandana Priya Srinivasa}
%% ============================================================
\textbf{Email:} \texttt{csrini10@asu.edu}.
\textbf{Role:} Side~B analysis.

\textbf{Charts.}
B2 (Black's weapons, grouped bars): curated the six classic Black
defenses (Sicilian, French, Caro-Kann, Scandinavian, Pirc, Modern)
and surfaced the Sicilian flip, the most-played opening in the
band where Black wins more than White.
B3 (Diverging extremes): specified the dual-CTE query for the top
five most-Black-favoring and top five most-White-favoring openings
($n \geq 500$), producing the visual swing of roughly 140pp.

%% ============================================================
\section{Parth Jain}
%% ============================================================
\textbf{Email:} \texttt{pjain121@asu.edu}.
\textbf{Role:} Side~B analysis.

\textbf{Charts.}
B4 (Bubble scatter, 60-opening spectrum): specified the
top-60-by-volume cutoff and the $x$=White~win\%, $y$=games,
size=games encoding, plus the click-to-pin interaction.
B6 (Factor variance, innovative): co-designed the project's
headline analytical chart, which compares the spread of White's
win\% across each candidate predictor. The output (color 3.9pp,
time control 0.4pp, Elo band 0.9pp, opening 19.1pp, rating diff
61.2pp) is the literal proof of the project's thesis.

%% ============================================================
\section{Joint Work and Effort}
%% ============================================================

The narrative pivots, the two innovative-chart concepts (A6
Tug-of-War and B6 Factor Variance), and the written report were
team-wide. Every chart was reviewed by at least one member outside
its owning sub-team before being considered final. Each member
contributed an estimated equal share of total project effort, with
the cross-cutting roles balancing the additional analytical
specification carried by the Side~B members on their charts.

\end{document}
